\item \model{OLMo 3}

\paragraph*{مقدمه و فلسفه بنیادین دادگان پیش‌آموزش (\en{Pre-Training Data Philosophy})}
در چارچوب چرخه کامل و شفاف مدل‌های زبانی سری \model{OLMo 3} (\en{Model Flow})، فاز پیش‌آموزش با هدف برپایی زیربنایی مستحکم جهت شکوفایی قابلیت‌های استدلال عمیق گام‌به‌گام (\en{Thinking}) و تعامل با ابزارها (\en{Tool Use}) در مراحل پس‌آموزش طراحی شده است \cite{olmo2024two}. به این منظور، پیکره عظیم \textbf{\en{Dolma 3}} در قالب یک استخر داده جامع به وسعت \textbf{۹.۳۱ تریلیون توکن} (شامل ۹.۹۷ میلیارد سند) گردآوری شده و ترکیب بهینه‌سازی‌شده نهایی آن موسوم به \textbf{\en{Dolma 3 Mix}} با حجم \textbf{۵.۹۳ تریلیون توکن} (شامل ۳.۸۷ میلیارد سند) مبنای پیش‌آموزش مدل‌های پایه در دو مقیاس ۷ و ۳۲ میلیارد پارامتر قرار گرفته است.

با توجه به این‌که بیش از ۹۰ درصد کل توان محاسباتی چرخه حیات مدل در فاز پیش‌آموزش مصرف می‌شود، استراتژی تدوین این پیکره بر دو اصل راهبردی استوار است:
\begin{itemize}
    \item \textbf{کفایت مقیاس داده‌ها (\en{Scale Sufficiency}):} یک منبع داده تنها در صورتی شایسته ورود به فاز پیش‌آموزش تلقی می‌شود که پتانسیل تولید حجم کافی از توکن‌ها را برای اثرگذاری در مقیاس تریلیونی دارا باشد؛ داده‌های کیفی ولی کم‌حجم به‌طور کامل برای مرحله میان‌آموزش (\en{Mid-training}) ذخیره می‌شوند.
    \item \textbf{پرهیز از داده‌های تکلیفی ساختاریافته (\en{Exclusion of Task-Structured Data}):} برای جلوگیری از مخدوش شدن نتایج ابلیشن و سوگیری‌های زودرس ارزیابی، داده‌های جفت پرسش‌وپاسخ (\en{QA}) و مکالمات دستوری از مرحله پیش‌آموزش حذف و به مراحل بعدی موکول شده‌اند.
\end{itemize}

\paragraph*{ترکیب‌بندی و مشخصات آماری پیکره پیش‌آموزش (\en{Dolma 3 Mix Composition})}
پیکره متنی پیش‌آموزش \en{Dolma 3 Mix} از شش منبع کلیدی به شرح جدول \ref{tab:olmo3_data_composition} استخراج، پالایش و متوازن شده است. در این فرآیند، استخر اولیه (\en{9T Pool}) با اعمال راهبردهای نمونه‌برداری به ترکیب نهایی (\en{6T Mix}) و همچنین یک نسخه سبک جهت پژوهش‌های مقیاس‌کوچک (\en{150B Mix}) تبدیل شده است.

\begin{table}[htbp]
\centering
\caption{ترکیب‌بندی و مشخصات آماری منابع مجموعه‌داده \en{Dolma 3} در استخر داده و ترکیب‌های آموزشی}
\label{tab:olmo3_data_composition}
\resizebox{\textwidth}{!}{%
\begin{tabular}{llcccccc}
\toprule
\textbf{منبع داده} & \textbf{نوع محتوا} & \multicolumn{2}{c}{\textbf{استخر داده (\en{9T Pool})}} & \multicolumn{2}{c}{\textbf{ترکیب نهایی (\en{6T Mix})}} & \multicolumn{2}{c}{\textbf{ترکیب آزمایشی (\en{150B Mix})}} \\
\cmidrule(lr){3-4} \cmidrule(lr){5-6} \cmidrule(lr){7-8}
& & \textbf{توکن} & \textbf{اسناد} & \textbf{توکن (درصد)} & \textbf{اسناد} & \textbf{توکن (درصد)} & \textbf{اسناد} \\
\midrule
\en{Common Crawl} & صفحات وب باز & 8.14T & 9.67B & 4.51T (76.1\%) & 3.15B & 121B (76.9\%) & 84.5M \\
\en{olmOCR Science PDFs} & اسناد علمی و مقالات & 972B & 101M & 805B (13.6\%) & 83.8M & 19.9B (12.6\%) & 2.25M \\
\en{StackEdu (Rebalanced)} & کدهای برنامه‌نویسی گیت‌هاب & 137B & 167M & 409B (6.89\%) & 526M & 11.1B (7.06\%) & 14.3M \\
\en{arXiv} & مقالات علمی ساختاریافته با لاتک & 21.4B & 3.95M & 50.8B (0.86\%) & 9.10M & 1.29B (0.82\%) & 247K \\
\en{FineMath 3+} & متون آموزشی ریاضیات در وب & 34.1B & 21.4M & 152B (2.56\%) & 95.5M & 4.10B (2.60\%) & 2.57M \\
\en{Wikipedia \& Wikibooks} & متون دانشنامه‌ای مرجع & 3.69B & 6.67M & 2.51B (0.04\%) & 4.24M & 64.6M (0.04\%) & 119K \\
\midrule
\textbf{مجموع کل (\en{Total})} & — & \textbf{9.31T} & \textbf{9.97B} & \textbf{5.93T (100\%)} & \textbf{3.87B} & \textbf{157B (100\%)} & \textbf{104M} \\
\bottomrule
\end{tabular}%
}
\end{table}

\paragraph*{خط لوله مهندسی و پالایش دادگان وب (\en{Web Data Pipeline})}
بخش اصلی استخر وب از ۱۰۴ دامپ (\en{Dump}) پایگاه \en{Common Crawl} (از سال ۲۰۱۳ تا پایان دسامبر ۲۰۲۴) گردآوری شده است. استخراج ساختار متنی و پیاده‌سازی متون عاری از برچسب‌های \en{HTML} با استفاده از کتابخانه \en{Resiliparse} \cite{bevendorff2018elastic} بر روی فایل‌های \en{WARC} انجام گرفت که به حجم خام اولیه \textbf{۲۵۲.۶ میلیارد سند} انجامید. 

فرایند فیلترینگ اکتشافی (\en{Heuristic Filtering}) با اقتباس و بهینه‌سازی خط لوله \en{DCLM} \cite{li2024dclm} در چند لایه پیاده‌سازی شد:
\begin{itemize}
    \item \textbf{فیلتر نشانی‌های وب (\en{URL Filtering}):} مسدودسازی دامنه‌ها و لینک‌های غیراخلاقی و بدافزار بر مبنای لیست سیاه مجموعه‌های \en{FineWeb} \cite{penedo2024fineweb} و \en{RefinedWeb} \cite{penedo2023refinedweb} (حذف حدود ۱٪ اسناد).
    \item \textbf{فیلترهای اکتشافی متن:} حذف اسناد نامتعارف از حیث طول (بسیار کوتاه یا بسیار بلند)، متون با نسبت پایین کاراکترهای الفبایی-عددی و محتواهای دچار تکرار مفرط درونی.
    \item \textbf{پالایش خطوط و عبارات کلیشه‌ای (\en{Line Modifiers}):} حذف سطرها و بندهای آلوده به هرزنامه‌های پرتکرار نظیر «\en{items in cart}» یا «\en{read more...}» و حذف اسنادی که بخش عمده ساختار خود را در این اصلاح از دست دادند.
    \item \textbf{جداسازی زبانی (\en{FastText Language Filter}):} به‌کارگیری مدل طبقه‌بند زبانی با آستانه اطمینان $0.65$ برای گزینش زبان انگلیسی.
    \item \textbf{قواعد کیفی \en{MADLAD-400}:} اعمال قواعد گزینش ۲ و ۵ برگرفته از پژوهش کودوگونتا و همکاران \cite{kudugunta2023madlad} جهت پالایش جملاتی با تعداد کلمات تمام‌بزرگ یا الگوهای منظم غیرعادی (\en{Cursed regex}).
\end{itemize}
این مراحل در مجموع با صرف ۱۰۳۰ ساعت محاسباتی روی ماشین‌های \en{i4i.32xlarge} ابری (معادل ۱۱,۳۰۰ دلار هزینه) منجر به نرخ حذف کلی ۸۴.۹٪ و بقای \textbf{۳۸.۷ میلیارد سند} گردید (جدول \ref{tab:olmo3_web_filtering}).

\begin{table}[htbp]
\centering
\caption{ریز تفکیک مراحل پالایش اکتشافی پیکره وب \en{Common Crawl}}
\label{tab:olmo3_web_filtering}
\begin{tabular}{lccc}
\toprule
\textbf{مرحله پالایش (\en{Pipeline Step})} & \textbf{اسناد حذف‌شده (میلیارد)} & \textbf{درصد حذف از کل استخر} & \textbf{درصد زمان پردازش} \\
\midrule
فیلترهای \en{URL} & 2.3 & 0.90\% & 1.68\% \\
فیلترهای طول سند & 103.4 & 40.42\% & 8.03\% \\
فیلترهای نماد و تراکم کاراکتر & 56.5 & 22.10\% & 4.13\% \\
تکرار درونی (\en{Internal Repetition}) & 32.1 & 12.53\% & 31.41\% \\
اصلاح‌کننده‌های خطوط (\en{Line Modifiers}) & 7.1 & 2.79\% & 10.00\% \\
فیلتر زبان انگلیسی (\en{FastText}) & 6.2 & 2.44\% & 14.30\% \\
فیلترهای \en{MADLAD-400} & 9.3 & 3.65\% & 5.87\% \\
طبقه‌بندهای کیفیت & 0.0 & 0.00\% & 24.58\% \\
\midrule
\textbf{مجموع کل (\en{Total})} & \textbf{213.9} & \textbf{84.9\% حذف (بقای 15.1\%)} & \textbf{100\% (1,030 ساعت)} \\
\bottomrule
\end{tabular}
\end{table}

\paragraph*{معماری سه‌سطحی حذف تکرار مقیاس‌بالا (\en{Three-Stage Deduplication})}
به منظور افزایش بهره‌وری مصرف توکن \cite{lee2022deduplicating}، حذف هم‌پوشانی‌های ضعیف به عنوان شاخص کیفیت نادقیق \cite{fang2025datasets}، و پیشگیری از افت کارایی ناشی از تکرار بیش از حد متون \cite{muennighoff2025scaling}، یک راهبرد سه‌سطحی با ابزار اختصاصی \textbf{\en{Duplodocus}} (توسعه‌یافته به زبان \en{Rust}) و \textbf{\en{bsade}} به کار رفت:
\begin{itemize}
    \item \textbf{حذف تکرار دقیق سراسری (\en{Exact Deduplication}):} در دو گام پی‌درپی با بهره‌گیری از هَش متنی ۱۲۸ بیتی؛ ابتدا درون هر یک از ۱۰۴ دامپ (حذف ۲۴٪ اسناد) و سپس به‌طور سراسری در میان کل دامپ‌ها (حذف ۴۳٪ اسناد باقی‌مانده). در نتیجه، ۶۶٪ از اسناد اولیه حذف و استخر به \textbf{۱۲.۷ میلیارد سند} کاهش یافت.
    \item \textbf{حذف تکرار فازی در سطح سند (\en{MinHash Fuzzy Deduplication}):} افراز داده‌ها به ۳۲ قطعه (\en{Shard})، توکنایزیشن با مدل \en{p50k} و استخراج ۵-گِرم‌ها. الگوریتم \en{MinHash LSH} با ۲۶ باند به اندازه ۱۱ برای هدف‌گذاری تشابه جاکارد $J \ge 0.80$ پیاده شد. در مرحله بعد، اعتبارسنجی دوبه‌دو برای کنترل مثبت کاذب بر پایه ۳-گِرم‌ها صورت پذیرفت؛ برای خوشه‌های ۵۰۰ سند یا بیشتر، ۲۰۰ باند به اندازه ۳۱ اعمال شد و برای خوشه‌های کوچک‌تر، فاصله جاکارد به صورت دقیق محاسبه گردید. با نگه‌داشتن جدیدترین نسخه سند بر مبنای تاریخ خزش، ۲۴٪ اسناد حذف و استخر به \textbf{۹.۸ میلیارد سند} رسید.
    \item \textbf{حذف زیررشته‌های تکراری با آرایه پسوندی (\en{Fuzzy Suffix Array Deduplication}):} با استفاده از \en{bsade} در ۵۶ بخش، تمام زیررشته‌های دارای طول ۵۰۰ بایت یا بیشتر که حداقل دو بار تکرار شده بودند نشانه‌گذاری شدند. طبق یک قاعده نوآورانه فازی، هرگاه یک بخش متنی از دو سو با توالی‌های تکراری ۵۰۰ بایتی محصور شده بود و بیش از ۸۰٪ آن بخش هم‌پوشانی تکراری داشت، کل توالی حذف شد تا قطعات ریز و ناقص از بین بروند. این گام موجب حذف ۱۴٪ از بایت‌های متنی شده و پیکره نهایی وب را به \textbf{۹.۷ میلیارد سند (۳۶.۵ ترابایت متن، معادل ۸ تریلیون توکن)} رساند.
\end{itemize}

\paragraph*{دسته‌بندی موضوعی، کیفی و ساختار ۴۸۰‌گانه سطل‌ها}
پس از حذف تکرار، استخر داده‌ها در یک ماتریس دوبعدی از موضوع و کیفیت افراز شد:
\begin{itemize}
    \item \textbf{طبقه‌بندی موضوعی (\en{Topic Classification}):} نگاشت متون به ۲۴ رده ساختاریافته برگرفته از \en{WebOrganizer} \cite{wettig2025organize}. جهت مقرون‌به‌صرفه بودن پردازش تریلیونی، یک طبقه‌بند \en{FastText} بر روی نمونه‌های نشاندار با \en{GPT-4.1} و \en{o4-mini} آموزش داده شد که به دقت و فراخوانی متوازن $0.762$ در جامعه آماری آزمون دست یافت.
    \item \textbf{ارزیابی کیفی (\en{Quality Tiering}):} با استفاده از مدل \en{FastText} آموزش‌دیده بر نمونه‌های مثبت شامل \en{OpenHermes-2.5} \cite{teknium2023openhermes}, \en{ELI5} \cite{fan2019eli5}, \en{UltraChat} \cite{ding2023enhancing} و \en{WildChat} \cite{zhao2024wildchat}، اسناد به ۲۰ بازه صدکی کیفی ۵ درصدی (\en{Vigintile}) تقسیم شدند.
\end{itemize}
تلاقی ۲۴ رده موضوعی و ۲۰ طبقه کیفی، \textbf{۴۸۰ سطل مجزا و متعامد} را شکل داد که دستمایه بهینه‌سازی توزیع در مرحله ترکیب گردید.

\paragraph*{پیکره مقالات و اسناد آکادمیک (\en{olmOCR Science PDFs})}
در یک پیشرفت شاخص نسبت به دادگان \en{peS2o} \cite{soldaini2023pes2o}، مدل \model{OLMo 3} از اسناد علمی پردازش‌شده توسط خط لوله نوین \textbf{\en{olmOCR}} \cite{poznanski2025olmocr} بهره می‌برد:
\begin{itemize}
    \item \textbf{استخراج متن با مدل‌های بینایی-زبانی:} از میان ۲۳۸ میلیون سند \en{PDF} آکادمیک اولیه (تا تاریخ دسامبر ۲۰۲۴)، اسناد با ابزار \en{olmOCR} (نسخه‌های ۰.۱.۴۹ تا ۰.۱.۵۳) تبدیل به متن خطی و تمیز شدند و کتابخانه \en{Poppler} فقط در موارد اضطراری به کار رفت. این گام ۱۶۰ میلیون سند تحویل داد.
    \item \textbf{حذف داده‌های هویتی بافت‌آگاه (\en{Context-Aware PII Removal}):} برای جداسازی داده‌های هویتی خصوصی از موارد عمومی، یک خط لوله مدل‌محور دو سطحی طراحی شد: مدل \en{Gemma-3-12B} صفحه اول سند را برای یافتن هویت‌های حساس پایش می‌کرد و مدل \en{Gemma-3-4B} با تحلیل ۵۰۰۰ کاراکتر ابتدایی، ماهیت سند را بازشناسی می‌نمود. این تفکیک هوشمند منجر به پالایش ۴.۹٪ اسناد و صیانت از حریم شخصی شد.
    \item \textbf{پالایش نهایی و پروانه نشر:} حذف اسناد با تراکم بیش از ۳۰٪ جدول، متون دارای بیش از ۲۰٪ محتوای صرفاً عددی، تبدیل جداول \en{Markdown} به کدهای ساختاریافته \en{HTML} و فیلترینگ سخت‌گیرانه برای حفظ اسناد دارای پروانه‌های آزاد و تجاری، مجموعه‌ای شامل \textbf{۱۰۸ میلیون سند علمی (۹۷۲ میلیارد توکن)} به بار آورد.
\end{itemize}

\paragraph*{صورت‌بندی ریاضی ترکیب داده‌ها و نمونه‌برداری آگاه از کیفیت (\en{Olmix \& Upsampling})}
انتخاب داده‌ها از میان ۴۸۰ سطل استخر وب و سایر منابع، از طریق حل یک مسئله بهینه‌سازی مقید ریاضی دو بخشی انجام پذیرفت:

\subparagraph*{۱. چارچوب ترکیب داده‌های مقید (\en{Olmix Framework}):}
با الهام از روش‌های ازدحامی \cite{liu2024regmix, ye2025data, diao2025climb}، یک کلونی از مدل‌های پروکسی ۳۰ میلیون پارامتری با بودجه محاسباتی بهینه چینچیلا ($5\times$ معادل ۳ میلیارد توکن) روی توزیع‌های دیریکله نمونه‌برداری‌شده از حوزه‌ها آموزش دیده و عملکرد آن‌ها بر حسب معیار بیت بر بایت (\en{Bits-per-Byte - BPB}) در آزمون \en{Base Easy Suite} اندازه‌گیری شد. سپس یک مدل خطی تعمیم‌یافته (\en{GLM}) برای هر تسک برازش گردید:
\begin{equation}
    \min_{\mathbf{w}} \frac{1}{|\mathcal{T}|} \sum_{t \in \mathcal{T}} g_t(\mathbf{w}) \quad \text{subject to} \quad \sum_{i} w_i = 1, \quad w_i \le \frac{M \cdot X_i}{Z_{\text{total}}}
\end{equation}
که در آن $\mathbf{w}$ بردار وزن دامنه‌ها، $g_t(\cdot)$ مدل برآورد \en{BPB} تسک $t$، $X_i$ تعداد کل توکن‌های در دسترس در دامنه $i$، $Z_{\text{total}} = 5.93\text{T}$ کل بودجه توکن و $M \in [4, 7]$ سقف مجاز فاکتور تکرار هر حوزه است. فرآیند در سه راند ترکیب شرطی (\en{Conditional Mixing}) تثبیت شد تا سهم وب به ۷۵٪ مقید شده و زبان‌های کدهای \en{StackEdu} \cite{allal2025smollm2, lozhkov2024starcoder2} و موضوعات مقالات علمی بهینه‌سازی شوند.

\subparagraph*{۲. نمونه‌برداری افزایشی بر مبنای توابع توانی-نمایی بریده‌شده (\en{Quality-Aware Upsampling}):}
به جای فیلترینگ ایستا (\en{Flat Filtering}) که اسناد را تنها به صورت دودویی حذف یا تک‌نسخه‌ای حفظ می‌کند، از خانواده توابع تحلیلی پیوسته و محدب زیر برای وزن‌دهی صدک‌های کیفی داده‌های وب استفاده شد:
\begin{equation}
    f_{p,\lambda}(x) = \begin{cases} 
        0, & x < a \\ 
        C(x - a)^p \cdot e^{\lambda(x - a)}, & x \ge a 
    \end{cases}
\end{equation}
به‌طوری‌که پارامترهای $p \ge 0$، $\lambda \ge 0$ و ضریب مقیاس $C$ از طریق حل عددی دستگاه معادلات مقید زیر استخراج می‌شوند:
\begin{equation}
    \int_{0}^{1} f_{p,\lambda}(x)\,dx = \frac{Z}{X}, \qquad \frac{1}{b}\int_{1-b}^{1} f_{p,\lambda}(x)\,dx \le M
\end{equation}
در این پیکربندی، آستانه $a = 0.40$ تعیین شد (دور ریختن ۴۰٪ داده‌های ضعیف)، $Z/X$ نسبت توکن تخصیص‌یافته توسط \en{Olmix}، و سقف تکرار برای $b$ درصد از مرغوب‌ترین داده‌های انتهایی برابر با $M = 7$ لحاظ گردید. در نهایت، فاکتور بازنمونه‌برداری هر سطل صدکی $[\,q^-,\, q^+\,]$ از رابطه زیر به دست آمد:
\begin{equation}
    \text{Upsampling Factor}_q = \frac{1}{q^+ - q^-} \int_{q^-}^{q^+} f_{p,\lambda}(x)\,dx
\end{equation}

\paragraph*{اعتبارسنجی تجربی و نتایج ابلیشن دادگان پیش‌آموزش}
صحت کارایی روش‌های فوق با آموزش مدل‌های ۱ میلیارد پارامتری روی ۱۰۰ میلیارد توکن در بنچمارک \en{Base Easy Suite} اعتبارسنجی شد (کاهش معیار \en{BPB} بیانگر بهبود کیفیت است). نتایج جدول \ref{tab:olmo3_pretrain_ablations} برتری قاطع رویکرد تلفیقی پیشنهادی را به اثبات می‌رساند.

\begin{table}[htbp]
\centering
\caption{نتایج مطالعات حذف (\en{Ablation}) در انتخاب استراتژی ترکیب داده و نمونه‌برداری کیفی}
\label{tab:olmo3_pretrain_ablations}
\begin{tabular}{lccc}
\toprule
\textbf{استراتژی پیکربندی داده} & \textbf{\en{QA Easy} (\en{BPB})} & \textbf{\en{Math Easy} (\en{BPB})} & \textbf{\en{Code Easy} (\en{BPB})} \\
\midrule
توزیع طبیعی داده‌ها (\en{Natural Distribution}) & 100.70 & 71.87 & 59.20 \\
ترکیب چندهدفه متوازن \en{Olmix} & \textbf{99.54} & \textbf{61.69} & \textbf{48.94} \\
\midrule
فیلترینگ ایستا ۵۰٪ برتر (تکرار ۱.۱ برابر) & 104.19 & 86.28 & 94.34 \\
فیلترینگ ایستا ۵٪ برتر (تکرار ۱۱.۱ برابر) & 106.53 & 84.31 & 93.02 \\
نمونه‌برداری کیفی پیوسته \model{OLMo 3} & \textbf{99.95} & \textbf{74.04} & \textbf{71.88} \\
\midrule
تلفیق با میانگین هندسی (\en{Geometric Mean}) & 100.42 & 78.18 & 81.34 \\
تلفیق با تابع نمایی ساده ($C e^{\lambda(x-a)}$) & 100.17 & 78.24 & 78.72 \\
\textbf{تلفیق توانی-نمایی بریده‌شده (\model{OLMo 3 Mix})} & \textbf{99.31} & \textbf{75.81} & \textbf{78.26} \\
\bottomrule
\end{tabular}
\end{table}

\paragraph*{پیکربندی محاسباتی و هایپرپارامترهای پیش‌آموزش}
آموزش مدل‌های پایه بر بستر زیرساخت پرسرعت \en{OLMo-core} در قالب مخرن توزیع‌پذیر \en{PyTorch FSDP2} اجرا گردید. مشخصات فنی این مرحله در جدول \ref{tab:olmo3_pretrain_hyperparameters} گزارش شده است. جهت ارزیابی پیوسته کیفیت مدل ۳۲ میلیارد پارامتری بدون نیاز به توقف‌های پرهزینه بازپخت، روش میانگین‌گیری دوره‌ای وزن‌ها از ۴ چک‌پوینت متوالی با فاصله ۱۰۰۰ گام \cite{li2025model} به کار گرفته شد.

\begin{table}[htbp]
\centering
\caption{هایپرپارامترهای فاز پیش‌آموزش مدل‌های پایه \model{OLMo 3}}
\label{tab:olmo3_pretrain_hyperparameters}
\begin{tabular}{lcc}
\toprule
\textbf{پارامتر تنظیماتی (\en{Hyperparameter})} & \textbf{\model{OLMo 3 Base-7B}} & \textbf{\model{OLMo 3 Base-32B}} \\
\midrule
تعداد لایه‌ها (\en{Layers}) & 32 & 64 \\
بعد پنهان ($d_{\text{model}}$) & 4096 & 5120 \\
سازوکار توجه (\en{Attention Heads}) & 32 سر Q / 32 سر KV & 40 سر Q / 8 سر KV (\en{GQA}) \\
طول پنجره زمینه (\en{Sequence Length}) & 8,192 توکن & 8,192 توکن \\
توجه پنجره لغزان (\en{Sliding Window Attention}) & پنجره ۴۰۹۶ در ۳/۴ لایه‌ها & پنجره ۴۰۹۶ در ۳/۴ لایه‌ها \\
کل توکن‌های آموزش فاز اول & 5.93 تریلیون توکن & 5.70 تریلیون توکن \\
برنامه نرخ یادگیری (\en{LR Schedule}) & کسینوسی تعدیل‌شده تا 5.9T & کسینوسی منقطع در 5.7T \\
اوج نرخ یادگیری (\en{Peak Learning Rate}) & $3.0 \times 10^{-4}$ & $6.0 \times 10^{-4}$ \\
حداقل نرخ یادگیری نهایی (\en{Final LR}) & $3.0 \times 10^{-5}$ & $6.21 \times 10^{-5}$ \\
اندازه دسته در سطح توکن (\en{Batch Size}) & 4,194,304 توکن & 8,388,608 توکن \\
دمای اوج آموزش ($\text{LR}^2 / \text{bsz}$) & $2.146 \times 10^{-14}$ & $4.292 \times 10^{-14}$ \\
ضریب \en{Z-Loss} & $10^{-5}$ & $10^{-5}$ \\
سخت‌افزار پردازشی & 1024 کارت \en{NVIDIA H100} & 1024 کارت \en{NVIDIA H100} \\
توان عملیاتی به ازای هر پردازنده & 7,700 توکن/ثانیه (43\% \en{MFU}) & 1,900 توکن/ثانیه (41\% \en{MFU}) \\
\bottomrule
\end{tabular}
\end{table}